\chapter{Medición de la radiación}

\section{Identificación y gestión de fuentes radiactivas}

Cabe señalar que las fuentes radiactivas se clasifican en fuentes encerradas o selladas y fuentes abiertas o no encerradas. Las primeras son aquellas en las que el material radiactivo está contenido dentro de una envoltura de suficiente resistencia mecánica para impedir que se establezca contacto con el radionúclido o que la sustancia radiactiva se disperse en las condiciones previsibles de utilización y desgaste. Una fuente abierta es toda fuente que no está sellada y que, en las condiciones normales de uso, puede producir contaminación; es decir, existe el riesgo de que parte del material radiactivo tenga contacto físico no deseado con cuerpos o materiales del medio que la rodea.

Respecto a los riesgos antes mencionados, puede decirse que una fuente encerrada, durante su uso normal, solo produce riesgo de irradiación externa mientras no pierda hermeticidad en su envoltura, mientras que una fuente abierta produciría riesgo de irradiación externa e interna, pues la posibilidad de contaminación y posterior ingestión o inhalación es mucho mayor.

Dentro de las finalidades de la medición de la radiación está la vigilancia individual, la cual se entiende como las actividades cuyo objetivo es estimar el equivalente de dosis y la incorporación de radionúclidos en las personas ocupacionalmente expuestas.

La vigilancia radiológica de los trabajadores ocupacionalmente expuestos a las radiaciones requiere de un sistema completo de medición, evaluación y registro de todas las exposiciones a las radiaciones que pueda sufrir cada individuo.

Una de las características que hace que la radiación sea peligrosa es que no es posible detectarla con nuestros cinco sentidos, por lo que requerimos de algún dispositivo especial o instrumento que lo haga por nosotros.

Para esto, se cuenta con varios tipos de detectores, que podríamos clasificar en:
\begin{itemize}
  \item \textbf{Dispositivos de vigilancia personal:} tienen como propósito llevar un registro de la dosis recibida por el trabajador o alertarlo en caso de que el nivel de radiación sea muy alto. Entre estos tenemos los dosímetros de bolsillo, dosímetros de película o placas monitoras.
  \item \textbf{Instrumentos de monitoreo:} se emplean como instrumentos de reconocimiento, a fin de evitar que los niveles de radiación sean mayores que los permitidos por la ley.
\end{itemize}

Para evaluar los equivalentes de dosis recibidos se emplean dosímetros individuales, los cuales deben ser llevados siempre por el personal ocupacionalmente expuesto mientras se encuentre trabajando.

\section{Dosímetros de bolsillo}

Los dosímetros de bolsillo son una pequeña cámara de ionización, tipo condensador, de forma cilíndrica, llena de aire y provista de un electrodo central.

Estos dispositivos deben cargarse antes de exponerlos a la radiación; esto se logra por medio de un cargador que les aplica un voltaje. Al exponer el dosímetro a la radiación, el aire de la cámara se ioniza y los iones se mueven hacia los electrodos, lo que ocasiona una pérdida de carga, la cual es proporcional a la cantidad de radiación recibida.

Los dosímetros de bolsillo pueden ser de lectura directa o de lectura indirecta. Los primeros están provistos de una fibra de cuarzo y de un pequeño microscopio, el cual permite observar el desplazamiento de la fibra a lo largo de una escala graduada. Este desplazamiento se debe a que, al perder carga, el electrodo atrae a la fibra de cuarzo; por lo tanto, el movimiento de la imagen de esta es una medida de la cantidad de radiación. En los de lectura indirecta, esta se efectúa mediante un dispositivo lector-cargador, el cual contiene un electrómetro que permite estimar la exposición o la dosis midiendo la carga perdida por el citado dosímetro. Estos dosímetros se utilizan para determinar dosis diarias.

Los dosímetros de bolsillo están sujetos a descargas cuando se tiran o golpean, por lo que debe tenerse cuidado en su manejo. Cuando los dosímetros son expuestos a la humedad, este aislamiento puede perderse, con el consiguiente aumento de la corriente de fuga. La rapidez de fuga normal en un buen dosímetro debe ser del orden de 3\% en un período de 48 horas.

\section{Dosímetros de película}

Los dosímetros de película consisten en un paquete de dos o tres placas fotográficas con plata (para rayos X, gamma y neutrones), protegidas de la luz y colocadas en un chasis. Estas películas constan de una base de acetato de celulosa o de algún plástico, revestida por ambos lados de una capa gelatinosa fotosensible (emulsión). De un lado se usa una emulsión muy sensible, mientras que del otro la emulsión es prácticamente insensible.

Este tipo de dosímetros se basa en el hecho de que la acción de la radiación ionizante o de sus emisiones secundarias produce un conglomerado microscópico de átomos de plata que constituye la imagen latente, la cual, al ser revelada, produce un oscurecimiento visible de la emulsión.

Una característica desfavorable de los dosímetros de película es que la dosis se puede determinar con un error de 20 a 30\%.

\section{Alarma sonora y comparación de dosímetros}

La alarma sonora tiene como característica anunciar en forma audible que se está en presencia de un campo de radiación alto. Estos monitores están calibrados de tal forma que, cuando están en presencia de un campo de radiación de 1 mR/h, emiten un ``bip'' cada 30 segundos.

Así, si un trabajador está en un campo de radiación de 2 mR/h, su alarma dará dos ``bips'' cada 30 segundos. En el caso de que un trabajador, sin saberlo, estuviese colocado en un campo de radiación alto, digamos 60 mR/h, su alarma estaría emitiendo 60 ``bips'' cada 30 segundos, es decir, dos cada segundo.

Estos monitores de alarma tienen dos posiciones. En la baja (que significa baja sensibilidad) se usan la mayor parte del tiempo. En la alta sensibilidad su comportamiento es similar, solo que emiten 30 ``bips'' en 30 segundos en un campo de 1 mR/h; esta escala se usa solo para inspeccionar paquetes.

En resumen, se pueden mencionar las siguientes ventajas y desventajas:

\begin{description}
  \item[Dosímetro de lectura directa.] \textbf{Ventaja:} información rápida y a la mano. \textbf{Desventajas:} si se golpea puede perder la información; si se recibe mucha dosis, se satura y pierde la lectura.
  \item[Dosímetro de película.] \textbf{Ventajas:} no se satura y no pierde la información al golpearse. \textbf{Desventajas:} necesita revelarse para ser leído y es dañado por la humedad.
  \item[Monitor de alarma.] \textbf{Ventaja:} alerta en forma audible si el nivel de radiación es alto. \textbf{Desventaja:} no proporciona una lectura del nivel de radiación.
\end{description}

\section{Detector Geiger--Müller}

Dentro de los instrumentos de monitoreo se encuentran el monitor con detector Geiger y el monitor con detector de centelleo.

El detector Geiger--Müller (GM) mide directamente el nivel de radiación al convertir la radiación ionizante en pulsos de corriente eléctrica y generalmente tiene tres escalas:
\[
0\text{--}10,\qquad 0\text{--}100,\qquad 0\text{--}1000\ \mathrm{mR/h}.
\]

Son aparatos confiables y, en el caso de la radiografía industrial, son indispensables para un trabajo seguro.

El botón de encendido tiene cinco posiciones: apagado, baterías (para comprobar la carga), $\times 100$, $\times 10$ y $\times 1$. La lectura de la aguja en la escala graduada de 0--10 se multiplica por 1, 10 o 100 según la posición en que se encuentre el selector.

Por ejemplo, si la aguja estuviera marcando 2 mR/h y el selector estuviera en la posición $\times 10$, significaría que la lectura correcta es de 20 mR/h.

\subsection{Ventajas y desventajas}

\textbf{Ventaja:} mide el nivel de radiación al cual se está exponiendo la persona.\\
\textbf{Desventaja:} no registra la dosis que ha recibido el individuo.

\section{Monitor con detector de centelleo}

Se basa en la detección por medio de los destellos luminosos producidos por la radiación nuclear en ciertos materiales. La ventaja de este tipo de instrumentos es que posee mayor eficiencia de conteo.
